% ECONOMICS_PROBLEM: {"id": "R13-342", "title": "Remote work and city economies", "jel_code": "R13", "source_id": "OP-0468"}
% !TeX program = xelatex
\documentclass[11pt,a4paper]{article}
\usepackage[margin=23mm]{geometry}
\usepackage{fontspec}
\setmainfont{Latin Modern Roman}
\usepackage{amsmath,amssymb,mathtools}
\usepackage{xcolor,enumitem,booktabs,longtable,array,xurl,hyperref}
\definecolor{muted}{HTML}{53636C}
\hypersetup{colorlinks=true,urlcolor=blue,linkcolor=blue}
\setlength{\parindent}{0pt}
\setlength{\parskip}{5pt}
\newcommand{\R}{\mathbb R}
\newcommand{\E}{\mathbb E}
\newcommand{\N}{\mathbb N}
\newcommand{\ind}{\mathbf 1}
\newcommand{\Var}{\operatorname{Var}}
\newcommand{\Cov}{\operatorname{Cov}}
\newcommand{\supp}{\operatorname{supp}}
\DeclareMathOperator*{\argmin}{arg\,min}
\DeclareMathOperator*{\argmax}{arg\,max}
\newcommand{\entrymeta}[3]{{\sffamily\small\color{muted}\textbf{\texttt{#1}}\quad|\quad #2\par #3\par}}
\newcommand{\Problemref}[1]{\texttt{#1}}
\newcommand{\JELref}[2]{\texttt{#2} (JEL \texttt{#1})}
\begin{document}
\section*{Remote work and city economies}
\textbf{Problem ID:} \texttt{R13-342}\quad \textbf{JEL:} \texttt{R13}\par
% BEGIN ECONOMICS STATEMENT
\label{problem:OP-0468}
{\sffamily\small\color{muted}Original subject group: Urban, housing and spatial economics\par}
\label{definitions:problem:30007750}
\textbf{Audit classification:} Background; proposed extension. Repaired profit/property double counting and stated baseline-cohort and terminal-value conventions.
\subsection*{Definitions and precise research target}
\textit{Editorial formalization.} Consider a metropolitan economy with occupations differing in remote-work feasibility. A specified permanent increase in permitted home-working days changes location, employment, commuting, office demand and local services over twenty years. For a fixed baseline household population including relevant owners and taxpayers, define monetary benefits \(b_i\) and weights \(\omega_i>0\) as follows. The proposed empirical target is \[\Delta W=E\sum_i\omega_i b_i(1).\] Identify or bound it using credible intervention variation and a declared model of price, selection and capacity responses. Financing, ownership income, resources and public services enter once. Local trial effects do not automatically identify full equilibrium responses; report selection, transport and distributional assumptions.

Define a fixed baseline household set including commuters, landlords, firm owners and local taxpayers, and an explicit external-area boundary; follow households who leave. Remote-work treatment specifies permitted days, job eligibility, enforcement and complementary equipment under each firm type. Expected twenty-year indirect utility includes wages, commuting/leisure, housing and ownership income, municipal services and taxes. Define \(b_i(1)\) by \(\mathcal U_i^0(b_i(1))=\mathcal U_i^1(0)\), requiring unique finite benefits, and \(\Delta W=E\sum_i\omega_i b_i(1)\). Include public resource costs through financing/services once and do not add separately counted profits or property appreciation. New remote permission can change work intensity and productivity; those are outcomes. A finite-period welfare target requires a terminal-value assumption if office conversion or capital persists beyond twenty years. All expectations use a stated baseline population and probability law. Identification means every admissible data-generating model with the same observed-data law yields the same displayed target; otherwise report the identified set. Exact equations, horizons and assumptions are editorial, rather than source-given canonical conjectures.
\subsection*{Variants and assumption changes}
\begin{enumerate}
\item \textbf{Editorial occupation variant.} Restrict permission by baseline remote feasibility and count noneligible service-worker effects. Post-policy occupation changes do not redefine the cohort.
\item \textbf{Editorial conversion variant.} Add feasible office-to-housing conversion with resource cost and zoning limits. Compare twenty-year welfare with fixed office stock.
\end{enumerate}
\subsection*{References and source audit}
\textbf{1.} Remote-work real-estate review exists; twenty-year metropolitan welfare intervention is editorial. Locator: NBER Working Paper 30662; abstract. Primary indexed metadata/abstract inspected; direct page failed. URL: \url{https://www.nber.org/papers/w30662}.
\begin{enumerate}\raggedright
\item\hypertarget{definitions:source:30007750:1}{} Stijn Van Nieuwerburgh. 2022. \emph{The Remote Work Revolution: Impact on Real Estate Values and the Urban Environment}. NBER Working Paper 30662; abstract.
\url{https://www.nber.org/papers/w30662}.
DOI: \url{https://doi.org/10.3386/w30662}.
\end{enumerate}

\subsection*{Common conventions: Source mathematical conventions}

These conventions apply to each entry unless explicitly replaced.

\textbf{Models and identification.} A primitive model \(m\) specifies all probability laws, feasible choices, information, equilibrium and measurement rules used by its target. The admissible class \(\mathcal M\) is fixed independently of outcomes. If \(P_O(m)\) is its observable probability law and \(\tau(m)\) the target, its identified set at observable law \(P\) is
\[
 \mathcal I_\tau(P)=\{\tau(m):m\in\mathcal M,\ P_O(m)=P\}.
\]
Point identification means that this set is a singleton. An empty set signals incompatibility of data and assumptions. Coordinate bounds need not describe the joint set. Bounds are sharp only if every retained target is attainable or arbitrarily approachable by compatible models. Finite bounds require explicit restrictions on unobserved quantities; unrestricted confounding, hidden wealth, extrapolation or arbitrary equilibrium selection can yield uninformative sets.

\textbf{Causal interventions.} A potential outcome \(Y_i(p)\) is the outcome under a complete policy regime \(p\), including timing, financing, information and market spillovers. The target population and units are fixed before treatment. With interference, use the complete assignment vector or a justified exposure mapping; individual no-interference notation alone cannot identify an economy-wide intervention. A difference of mean potential outcomes does not require the cross-world joint distribution. Conditioning on post-treatment migration, survival, enrollment or employment changes the estimand and needs separate justification.

\textbf{Statistical statements.} An estimator, confidence set, forecast or test specifies a sampling law, model class and loss/coverage criterion. Uniform \((1-\alpha)\)-coverage means \(\inf_{m\in\mathcal M}\Pr_m\{\tau(m)\in C_n\}\geq1-\alpha\), or an explicitly asymptotic version. The whole parameter space trivially has coverage, so informativeness requires a stated width, risk or power objective. A forecasting improvement is relative to a named benchmark, information set and evaluation population.

\textbf{Equilibrium and optimization.} An equilibrium comprises feasible plans, optimality, beliefs where needed, budget constraints and all market/resource clearing. The primitives or a stated benchmark define each correspondence \(\mathcal E(p;m)\). Nonemptiness is assumed or proved, never inferred just from compact policy sets. A selected-equilibrium target uses a declared selection rule; a robust target quantifies over all permitted equilibria. Use suprema or infima unless attainment is established. An \(\varepsilon\)-optimal policy is feasible and within \(\varepsilon\) of the finite optimum value. Report an empty feasible set separately.

\textbf{Welfare and accounting.} Normative utility, interpersonal normalization, social weights, numeraire, discounting and the use of public receipts are primitives. Behavioral utility need not coincide with the welfare criterion. Household welfare includes ownership income and rents once; adding profits again to compensating variation generally double-counts them. A monetary equivalent is defined by an explicit transfer and price convention, with monotonicity and range conditions. Average effects, mechanism contributions, transfers and welfare losses are different targets. Infinite sums require convergence and dynamic feasible plans require terminal or no-Ponzi conditions.

\textbf{Source versions.} A working-paper date, revision date and journal publication year are distinguished. A DOI for a journal version is not automatically the DOI of an earlier manuscript. Locators refer to the stated version and printed pages where specified. A metadata-only check does not verify a claim about a full-text conclusion. Every entry's source subsection narrows the provenance claim accordingly.
% END ECONOMICS STATEMENT
\end{document}
